\section{Conditioned physical paths and their Gaussian bridges}
\label{sec:stationary-window-bridges}
We identify the common conditional law of
Theorem~\ref{thm:stationary-physical-windows}. The conditioning remains
an exact observation at deterministic physical time, with the same
shrinking standardized endpoint box. The result is a functional limit,
not merely a comparison of two conditional measures. Its proof pays
for the rare denominator in the tightness and clock estimates.

For a positive definite $3\times3$ matrix $C$ and $\xi\in\R^3$, let
\begin{equation}\label{eq:Gaussian-bridge-definition}
 \mathbb B_{C,\xi}(s)=s\xi+C^{1/2}(W(s)-sW(1)),\qquad 0\le s\le1,
\end{equation}
where $W$ is standard three-dimensional Brownian motion. This law on
$\mathcal C=C([0,1],\R^3)$ has mean $s\xi$ and covariance
$(\min\{s,t\}-st)C$. We use the supremum norm on $\mathcal C$
and the bounded-Lipschitz distance on probability laws.

\subsection{Tightness after conditioning on the endpoint}
Under the hypotheses of
Proposition~\ref{prop:window-conditional-characteristic}, let $X_m$
be the polygonal interpolation with
\[
 X_m(j/m)=P_m S_jh/\sqrt m,\qquad 0\le j\le m.
\]
Retain $E_m$, $C_m$ and $\varepsilon_m$ from that proposition.

\begin{lemma}[A trigonometric tail test]
\label{lem:bridge-trigonometric-tail}
For each positive integer $p$ there are $c_p,C_p>0$ such that
\[
 c_p\one_{\{|x|\ge1\}}
 \le F_p(x):=\int_0^1(1-\cos(ux))^p\dd u
 \le C_p|x|^{2p}.
\]
The integrand is a trigonometric polynomial of degree $p$ in $ux$.
\end{lemma}
\begin{proof}
The upper bound follows from $1-\cos v\le v^2/2$.
The function $F_p$ is continuous, even, and strictly positive away
from zero. On $1\le|x|\le4\pi$ it has a positive minimum.
For $|x|\ge4\pi$, write it as the average of the nonnegative,
$2\pi$-periodic function $(1-\cos v)^p$ over an interval of length
$|x|$. The complete periods occupy at least half that length and
have a positive mean. This gives the lower bound. Expanding the
power gives the last assertion.
\end{proof}

\begin{theorem}[A collision bridge under an unsmoothed endpoint window]
\label{thm:collision-window-bridge}
Uniformly in $R$, the initial densities \eqref{eq:bridge-initial-class},
the matrices and windows in
Proposition~\ref{prop:window-conditional-characteristic},
\begin{equation}\label{eq:collision-window-bridge}
 d_{\rm BL}\left(\operatorname{Law}_{\sigma\nu}(X_m\mid E_m),
                         \operatorname{Law}(\mathbb B_{C_m,\xi})\right)
                                      \longrightarrow0.
\end{equation}
No independence of the actual collision increments is assumed.
\end{theorem}
\begin{proof}
For a fixed finite list of times, use
\eqref{eq:window-conditional-characteristic} with bounded test
frequencies and the corresponding consecutive blocks. Rounding times
to collision indices and polygonal interpolation change a fixed
finite list by $O(m^{-1/2})$, since $h_R$ and $P_m$ are bounded.
The displayed Gaussian increment formula proves the finite-dimensional
bridge limit, uniformly in the parameters.

For tightness take any collision-grid increment of length
$\Delta=(k-j)/m$. Apply Lemma~\ref{lem:bridge-trigonometric-tail}
with $p=4$ to each of its three coordinates divided by $a/\sqrt3$.
The characteristic functions involved have three consecutive blocks
and test frequencies at most $4\sqrt3/a$. Therefore, whenever
$4\sqrt3/a\le c(C_0)m^{1/200}$,
\begin{equation}\label{eq:conditional-increment-tail}
 \mathbb P_{\sigma\nu}\{|X_m(k/m)-X_m(j/m)|>a\mid E_m\}
          \le C_A a^{-8}\Delta^4+C_A\varepsilon_m.
\end{equation}
Indeed the corresponding Gaussian bridge increment has mean
$\Delta\xi$, covariance $\Delta(1-\Delta)C_m$, and eighth moment
at most $C_A\Delta^4$. The finite trigonometric coefficient sum is
fixed, so its Fourier comparison costs $C_A\varepsilon_m$,
not a negative power of the endpoint probability.

Set
\begin{equation}\label{eq:bridge-tightness-scales}
 \eta=\frac1{200},\qquad \alpha=\frac18,\qquad
 J_m=\lfloor\eta\log_2m\rfloor.
\end{equation}
Use the nested grids with vertices
$\lfloor ml2^{-j}\rfloor/m$, $0\le l\le2^j$, for
$j\le J_m$. Each adjacent increment has length at most $2^{1-j}$.
Fix $\epsilon>0$ and assign threshold
$a_j=\epsilon2^{-\alpha j}$. For all sufficiently large $m$, all
frequencies in \eqref{eq:conditional-increment-tail} are permitted,
because their largest order is
$C_\epsilon m^{\eta\alpha}=C_\epsilon m^{1/1600}$,
strictly below $m^{1/200}$. A union bound over levels
$j_0\le j\le J_m$ and their adjacent increments gives
\begin{equation}\label{eq:bridge-coarse-grid-bound}
 C_A\epsilon^{-8}\sum_{j\ge j_0}2^{-2j}
       +C_A2^{J_m}\varepsilon_m
 \le C_A\epsilon^{-8}2^{-2j_0}
       +C_A m^{-9/2800}\sqrt{\log(2+m)}.
\end{equation}
Here $4(1-2\alpha)-1=2$ and $23/2800-1/200=9/2800$.
On the complementary event, dyadic chaining bounds the modulus of
the polygon through the finest grid vertices by a constant times
$\epsilon2^{-\alpha j_0}$ on intervals of length $2^{-j_0}$:
pass from each vertex to its ancestor at level $j_0$ and sum at most
two increments per finer level. Adjacent ancestors add at most a
fixed number of level-$j_0$ increments. Rounding the nested grids
only changes these constants.

It remains to control the motion inside the finest cells. This is
where ordinary tightness without a rare-event budget would be
insufficient. Let $\widetilde X_m$ be the polygon through those
vertices. Lemma~\ref{lem:fixed-even-maximal}, at the fixed order
$128$, gives on each cell, under the bounded density $\sigma$,
\[
 \mathbb E\sup_{s\text{ in the cell}}
          |X_m(s)-\widetilde X_m(s)|^{128}
                                     \le C2^{-64J_m}.
\]
The bound follows from deterministic-window partial sums and their
endpoint increment; the largest cell has $O(m2^{-J_m})$ collisions.
There are $2^{J_m}$ cells. Since
$\mathbb P_{\sigma\nu}(E_m)\ge c_A m^{-6/25}$,
\begin{equation}\label{eq:bridge-fine-grid-bound}
 \mathbb P_{\sigma\nu}
 \{\|X_m-\widetilde X_m\|_\infty>\epsilon\mid E_m\}
 \le C_A\epsilon^{-128}m^{6/25}2^{-63J_m}
 \le C_A\epsilon^{-128}m^{-3/40}.
\end{equation}
The strict exponent is $63/200-6/25=3/40$.
First let $m\to\infty$ in \eqref{eq:bridge-coarse-grid-bound}
and \eqref{eq:bridge-fine-grid-bound}, and then let $j_0\to\infty$.
They prove conditional tightness in $\mathcal C$.

Finally suppose uniformity in \eqref{eq:collision-window-bridge}
failed. Along a violating sequence select a subsequence for which
$C_m$ and $\xi$ converge; these lie in compact sets of uniformly
positive matrices and bounded vectors. Tightness and the uniform
finite-dimensional estimate identify every subsequential law as the
bridge with those limiting parameters. The comparison bridge laws
also converge: couple them by the same $W$ in
\eqref{eq:Gaussian-bridge-definition} and use continuity of the
positive matrix square root. This contradicts the violating
bounded-Lipschitz distance and completes the proof.
\end{proof}

\subsection{The whole physical clock on the same stationary space}
Use the stationary suspension $(\mathcal X_R,\mathbb P_R)$ and the
four exact endpoint events $A_t^j$ of
\eqref{eq:stationary-four-events}, with half-width $t^{-2/25}$.
Let $\mathsf X_{t,R}$ be the continuous polygon through the values
\[
 t^{-1/2}V_R(k),\qquad k=0,1,\ldots,\lfloor t\rfloor,
 \quad\hbox{and }t^{-1/2}V_R(t),
\]
at times $k/t$ and $1$, where
$V_R(u)=(W_R(u),C_R(u)-u/\bar\tau_R)$.
The last point is included only once when $t$ is an integer.
Its distance from the original physical step path is $O(t^{-1/2})$:
the minimum flight length is positive and a unit physical time interval
has bounded displacement and bounded collision count.

Let $m=m_R(t)$ and let $\mathsf Y_{t,R}$ be the collision polygon
with values $t^{-1/2}B_RS_jh_R(x)$ at times $j/m$.
The state $x$ is the actual outgoing collision preceding the stationary
observation, not an independently chosen initial state.

\begin{lemma}[A uniform pathwise physical-to-collision comparison]
\label{lem:whole-physical-clock-coupling}
For every fixed $P>0$,
\begin{equation}\label{eq:whole-physical-clock-coupling}
 \mathbb P_R\{\|\mathsf X_{t,R}-\mathsf Y_{t,R}\|_\infty
                                  >C_Pt^{-1/8}\}\le C_Pt^{-P}.
\end{equation}
Consequently the same exceptional probability under any of the four
endpoint conditionings is at most $C_{A,P}t^{-P+6/25}$.
\end{lemma}
\begin{proof}
Use Proposition~\ref{prop:stationary-observation-coupling} at each
integer observation length $k\ge t^{1/4}$ and at $t$. In unnormalized
variables its physical-to-collision error is $C_Q k^{3/8}$ outside
probability $C_Q k^{-Q}$. Summing over at most $t+2$ observation
lengths gives $C_Qt^{1-Q/4}$. Take $Q>4(P+1)$.
Every retained error is at most $C_Qt^{3/8}$.

For $k<t^{1/4}$ both the physical vector and its deterministic collision
reference have size at most $C(1+k)$: use bounded speed and cell-position
terms, the positive minimum flight length, and boundedness of $h_R$.
Their difference is therefore at most $Ct^{1/4}$, which is smaller
than the same threshold. At the physical grid point $k/t$ the index
$\lfloor mk/t\rfloor$ differs from $m_R(k)$ by at most two.
This changes the collision sum by a bounded amount. Between adjacent
physical grid points, the physical interpolation and the collision
polygon each change by a bounded amount before normalization.
Division by $\sqrt t$ proves \eqref{eq:whole-physical-clock-coupling}.
No union over a continuum of clocks has been taken. The lower bounds
in \eqref{eq:stationary-window-denominator} prove the last assertion.
\end{proof}

\begin{theorem}[Stationary physical bridges in exact endpoint windows]
\label{thm:stationary-physical-bridge}
For every $A<\infty$, uniformly for $R\in I$, $|\xi|\le A$, and
$j\in\{{\rm coll},{\rm ret},{\rm last},{\rm phys}\}$,
\begin{equation}\label{eq:stationary-physical-bridge}
 d_{\rm BL}\left(\operatorname{Law}_{\mathbb P_R}
          (\mathsf X_{t,R}\mid A_t^j),
               \operatorname{Law}(\mathbb B_{\mathcal V_R,\xi})\right)
                                      \longrightarrow0.
\end{equation}
In particular the actual physical observation conditions the entire
physical displacement and collision-fluctuation path to a Gaussian
bridge with terminal value $\xi$.
\end{theorem}
\begin{proof}
For $A_t^{\rm coll}$ apply
Theorem~\ref{thm:collision-window-bridge} under the exact collision
marginal $\sigma_R\nu$, $\sigma_R=\tau_R/\bar\tau_R$, with
$P_m=\sqrt{m/t}\,B_R$. This matrix has the uniformly bounded
augmentation used in Theorem~\ref{thm:stationary-physical-windows}.
Its covariance is
$C_m=(m/t)B_R\Gamma_RB_R^{\mathsf T}=\mathcal V_R+O(t^{-1})$,
and its endpoint window widths are comparable to $m^{-2/25}$.
The resulting polygon is exactly $\mathsf Y_{t,R}$.
Lemma~\ref{lem:whole-physical-clock-coupling}, with $P=2$, replaces
it by $\mathsf X_{t,R}$ under this conditioning, since the conditional
exceptional probability is $O(t^{-2+6/25})$. Continuity of Gaussian
bridge laws in their positive covariance gives the stated limit.

Theorem~\ref{thm:stationary-physical-windows} and the triangle
inequality bound the other three conditional initial laws against the
collision-conditioned law in total variation by
\[
 C_A t^{-23/2800}\sqrt{\log(2+t)}.
\]
Pushforward by the same physical path cannot increase this distance. This proves all four assertions. The return-conditioned
statistic remains the actual physical path on the original stationary
suspension; neither its initial section point nor its elapsed age is
resampled.
\end{proof}

\subsection{Actual denominators for a class of path selectors}
\begin{corollary}[Path-weighted window denominators]
\label{cor:bridge-path-weighted-denominator}
Let $F:\mathcal C\to[0,\infty)$ be fixed, bounded and continuous.
Uniformly on the parameter sets of the preceding theorem,
\begin{align}
 \mathbb E_R[F(\mathsf X_{t,R})\one_{A_t^j}]
  =8t^{-6/25}g_{\mathcal V_R}(\xi)
        \left(\mathbb E F(\mathbb B_{\mathcal V_R,\xi})+o(1)\right).
 \label{eq:bridge-path-weighted-denominator}
\end{align}
Whenever the bridge expectation has a positive uniform lower bound,
the conditional path law with weight $F(\mathsf X_{t,R})\one_{A_t^j}$
converges to the bridge law weighted by $F$. The relative comparison
between any two of these weighted endpoint events retains the rate
$O_A(t^{-23/2800}\sqrt{\log(2+t)})$, with a constant depending on
the stated lower bound and $\|F\|_\infty$.
\end{corollary}
\begin{proof}
Uniform weak convergence and uniform tightness give convergence of
expectations of every fixed bounded continuous function: on a common
compact path set approximate it uniformly by bounded Lipschitz
functions, and use boundedness outside that set. Multiply this
conclusion by \eqref{eq:stationary-window-denominator}.
For the reweighted law apply it also to $FG$ for a bounded continuous
path test $G$, and divide by the positive denominator.
For the comparison, the difference of the two unnormalized densities
is supported on the original symmetric difference and is at most
$\|F\|_\infty$ there. Normalization costs at most twice its integral
divided by either weighted mass. The relative event theorem and
\eqref{eq:bridge-path-weighted-denominator} give the claimed rate.
\end{proof}

\begin{corollary}[A concrete nonsingle-time indicator class]
\label{cor:bridge-cylinder-events}
Fix $0<s_1<\cdots<s_r<1$ and bounded boxes
$D_1,\ldots,D_r\subset\R^3$ with nonempty interiors. Define the
exact path event
\[
 D_t=\bigcap_{l=1}^r\{t^{-1/2}V_R(ts_l)\in D_l\}.
\]
Then, uniformly for $R\in I$, $|\xi|\le A$, and all four endpoint
events,
\begin{equation}\label{eq:bridge-cylinder-denominator}
 \mathbb P_R(D_t\cap A_t^j)
 =8t^{-6/25}g_{\mathcal V_R}(\xi)
       \mathfrak b_R(\xi;\mathbf s,\mathbf D)(1+o(1)),
\end{equation}
where
\[
 \mathfrak b_R(\xi;\mathbf s,\mathbf D)
  =\int_{\prod_lD_l}
      g_{\Sigma_R}\bigl((y_l-s_l\xi)_{l=1}^r\bigr)
                           \prod_l\dd y_l,\quad
 (\Sigma_R)_{lk}=(\min\{s_l,s_k\}-s_ls_k)\mathcal V_R.
\]
The factor $\mathfrak b_R$ has a positive uniform lower bound.
The conditional physical path converges to the bridge restricted to
these cylinder boxes. The conditional initial laws for
$D_t\cap A_t^j$ and $D_t\cap A_t^{j'}$ are within
$O_A(t^{-23/2800}\sqrt{\log(2+t)})$ in total variation.
\end{corollary}
\begin{proof}
For distinct interior times the scalar bridge covariance matrix is
positive definite. This follows, for example, by expressing its values
as independent Gaussian increments conditioned on their total: no
nonzero linear combination of the interior partial sums is a multiple
of that total. Uniform ellipticity of $\mathcal V_R$ gives positive
definiteness of $\Sigma_R$. The Gaussian density is uniformly positive
on any fixed compact subbox of $\prod_l\operatorname{int}D_l$, for
$R\in I$ and $|\xi|\le A$. This proves the lower bound.
The cylinder boundary has zero measure under every limiting bridge.
The actual observations differ from the interpolated values by
$O(t^{-1/2})$ uniformly. Enlarging and contracting the fixed cylinder
boxes therefore shows that the path convergence applies to the exact
physical indicator and to that indicator times a bounded continuous
path test. Uniformity follows by
the same compact-parameter subsequence argument; every subsequential
Gaussian has this boundary-null property. These observations prove
the denominator and the restricted bridge limit. Finally
$(D_t\cap A_t^j)\mathbin\triangle(D_t\cap A_t^{j'})
\subset A_t^j\mathbin\triangle A_t^{j'}$; divide the inherited
relative bound by the newly proved positive cylinder factor.
\end{proof}

\paragraph{Relation to the raw conditioning problem.}
The result supplies an actual conditional functional limit and explicit
weighted denominators for fixed interior cylinder observations, without
assuming their long pullbacks have controlled variation. The endpoint
width remains $t^{21/50}$ in physical coordinates. Fixed lattice labels,
microscopic time windows, the pointwise common signed raw correction,
and the complete fixed-return Fourier complement still require their
own estimates. No restriction on those original target events has been
silently imposed: the theorem above concerns the displayed window class
and leaves the original raw mixed-density endpoint intact.
